\documentclass[10pt,a4paper]{amsart}
\usepackage{amsmath,amssymb,amsthm,graphicx}
% Retained preamble for the Stacks Project chapter {CH} (shared project preamble).
% The chapter file pinned-original-source.tex only carries \input{{preamble}} on its first line,
% so no preamble text lives inside the chapter itself.  The Stacks Project uses one single shared
% preamble for every chapter: the snapshot stages it three times (sources/stacks-chow/preamble.tex,
% sources/stacks-etale-cohomology/preamble.tex and focused-seeds/AGP/sources/stacks-snapshot-20261011/
% preamble.tex) and those copies are identical in size.  The block below is copied mechanically,
% line for line, from sources/stacks-chow/preamble.tex lines 14-18 and 154-217: the Xy-pic setup and
% the native theorem declarations together with the mathematical macros.  Nothing was added, removed
% or redefined; the wrapper class and the amsmath/amssymb/amsthm/graphicx base packages are supplied
% by the assembler.
% For commutative diagrams we use Xy-pic
% Chapter: stacks-cohomology.tex -- Cohomology of Sheaves.
\usepackage[all]{xy}

% We use 2cell for 2-commutative diagrams.
\xyoption{2cell}

% Theorem environments.
%
\theoremstyle{plain}
\newtheorem{theorem}[subsection]{Theorem}
\newtheorem{proposition}[subsection]{Proposition}
\newtheorem{lemma}[subsection]{Lemma}

\theoremstyle{definition}
\newtheorem{definition}[subsection]{Definition}
\newtheorem{example}[subsection]{Example}
\newtheorem{exercise}[subsection]{Exercise}
\newtheorem{situation}[subsection]{Situation}

\theoremstyle{remark}
\newtheorem{remark}[subsection]{Remark}
\newtheorem{remarks}[subsection]{Remarks}

\numberwithin{equation}{subsection}

% Macros
%
\def\lim{\mathop{\mathrm{lim}}\nolimits}
\def\colim{\mathop{\mathrm{colim}}\nolimits}
\def\Spec{\mathop{\mathrm{Spec}}}
\def\Hom{\mathop{\mathrm{Hom}}\nolimits}
\def\Ext{\mathop{\mathrm{Ext}}\nolimits}
\def\SheafHom{\mathop{\mathcal{H}\!\mathit{om}}\nolimits}
\def\SheafExt{\mathop{\mathcal{E}\!\mathit{xt}}\nolimits}
\def\Sch{\mathit{Sch}}
\def\Mor{\mathop{\mathrm{Mor}}\nolimits}
\def\Ob{\mathop{\mathrm{Ob}}\nolimits}
\def\Sh{\mathop{\mathit{Sh}}\nolimits}
\def\NL{\mathop{N\!L}\nolimits}
\def\CH{\mathop{\mathrm{CH}}\nolimits}
\def\proetale{{pro\text{-}\acute{e}tale}}
\def\etale{{\acute{e}tale}}
\def\QCoh{\mathit{QCoh}}
\def\Ker{\mathop{\mathrm{Ker}}}
\def\Im{\mathop{\mathrm{Im}}}
\def\Coker{\mathop{\mathrm{Coker}}}
\def\Coim{\mathop{\mathrm{Coim}}}

% Boxtimes
%
\DeclareMathSymbol{\boxtimes}{\mathbin}{AMSa}{"02}

%
% Macros for moduli stacks/spaces
%
\def\QCohstack{\mathcal{QC}\!\mathit{oh}}
\def\Cohstack{\mathcal{C}\!\mathit{oh}}
\def\Spacesstack{\mathcal{S}\!\mathit{paces}}
\def\Quotfunctor{\mathrm{Quot}}
\def\Hilbfunctor{\mathrm{Hilb}}
\def\Curvesstack{\mathcal{C}\!\mathit{urves}}
\def\Polarizedstack{\mathcal{P}\!\mathit{olarized}}
\def\Complexesstack{\mathcal{C}\!\mathit{omplexes}}
% \Pic is the operator that assigns to X its picard group, usage \Pic(X)
% \Picardstack_{X/B} denotes the Picard stack of X over B
% \Picardfunctor_{X/B} denotes the Picard functor of X over B
\def\Pic{\mathop{\mathrm{Pic}}\nolimits}
\def\Picardstack{\mathcal{P}\!\mathit{ic}}
\def\Picardfunctor{\mathrm{Pic}}
\def\Deformationcategory{\mathcal{D}\!\mathit{ef}}
\usepackage{cleveref}
\crefname{originalreference}{original reference}{original references}
\Crefname{originalreference}{Original reference}{Original references}
\makeatletter
\newlabel{situation-eta}{{0GT4}{source}{}{sourcepointer-1}{}}
\newlabel{situation-eta@cref}{{[originalreference][0][]0GT4}{[1][0][]source}{}{}{}}
\makeatother
\crefname{originalreference}{original reference}{original references}
\Crefname{originalreference}{Original reference}{Original references}
\makeatletter
\newlabel{derived-lemma-localization-subcategory}{{0GSL}{source}{}{sourcepointer-2}{}}
\newlabel{derived-lemma-localization-subcategory@cref}{{[originalreference][0][]0GSL}{[1][0][]source}{}{}{}}
\makeatother
\crefname{originallemma}{original lemma}{original lemmas}
\Crefname{originallemma}{Original lemma}{Original lemmas}
\makeatletter
\newlabel{lemma-K-flat-resolution}{{06YF}{source}{}{sourcepointer-3}{}}
\newlabel{lemma-K-flat-resolution@cref}{{[originallemma][0][]06YF}{[1][0][]source}{}{}{}}
\makeatother
\crefname{originallemma}{original lemma}{original lemmas}
\Crefname{originallemma}{Original lemma}{Original lemmas}
\makeatletter
\newlabel{lemma-I-torsion-free}{{0GT5}{source}{}{sourcepointer-4}{}}
\newlabel{lemma-I-torsion-free@cref}{{[originallemma][0][]0GT5}{[1][0][]source}{}{}{}}
\makeatother
\crefname{originallemma}{original lemma}{original lemmas}
\Crefname{originallemma}{Original lemma}{Original lemmas}
\makeatletter
\newlabel{lemma-eta-qis}{{0F8P}{source}{}{sourcepointer-5}{}}
\newlabel{lemma-eta-qis@cref}{{[originallemma][0][]0F8P}{[1][0][]source}{}{}{}}
\makeatother
\crefname{originalreference}{original reference}{original references}
\Crefname{originalreference}{Original reference}{Original references}
\makeatletter
\newlabel{categories-lemma-properties-left-localization}{{04VG}{source}{}{sourcepointer-6}{}}
\newlabel{categories-lemma-properties-left-localization@cref}{{[originalreference][0][]04VG}{[1][0][]source}{}{}{}}
\makeatother
\title{Existence of the derived functor L-eta\_I computed by the eta construction on complexes of I-torsion free modules}
\author{The Stacks Project Authors}
\date{stacks-project:cohomology.tex:local-snapshot-20261011; GNU FDL 1.2 or later}
\begin{document}
\maketitle
\noindent\textit{Editorial note.} All delivered mathematical text is the author's own, copied verbatim from the pinned Stacks Project chapter file cohomology.tex; no statement or proof step was rewritten, reordered or removed.\par\medskip
\noindent\textit{Editorial note.} The counted span contains no commutative diagram and cites only cross-chapter results, each recorded in assembly.external\_references with its canonical tag; no stub replaces a load-bearing argument.\par\medskip
\noindent\textit{Editorial note.} The proof cites no bibliography item inside the counted span.\par\medskip
\noindent\textit{Allowed standard prerequisites (editorial).}\begin{itemize}
\item Localization of a full triangulated subcategory at quasi-isomorphisms. Stacks Project, Derived Categories, Lemma derived-lemma-localization-subcategory (tag 0GSL). The criterion that reduces the equivalence S\textasciicircum{}\{-1\}K(T) = D(O\_X) to the existence of torsion-free resolutions.
\item Existence of K-flat resolutions. Stacks Project, Cohomology of Sheaves, Lemma lemma-K-flat-resolution (tag 06YF). Provides the flat (hence I-torsion free) complex quasi-isomorphic to a given complex.
\item Characterizations of I-torsion freeness and quasi-isomorphism invariance of eta. Stacks Project, Cohomology of Sheaves, Lemmas lemma-I-torsion-free (tag 0GT5) and lemma-eta-qis (tag 0F8P). The third characterization links flatness to torsion freeness, and eta-qis is what makes the functor descend to the derived category.
\end{itemize}\medskip
\ifdefined\hypertarget\hypertarget{sourcepointer-1}{}\fi
\noindent\textit{Source reference.} Original reference 0GT4: sources/stacks-chow/tags.txt, source lines 20954--20954. Cross-chapter standard prerequisite cited by the author inside the counted span; it is used as a pointer and its own proof is not part of this item.\par\medskip
\ifdefined\hypertarget\hypertarget{sourcepointer-2}{}\fi
\noindent\textit{Source reference.} Original reference 0GSL: sources/stacks-chow/tags.txt, source lines 20936--20936. Cross-chapter standard prerequisite cited by the author inside the counted span; it is used as a pointer and its own proof is not part of this item.\par\medskip
\ifdefined\hypertarget\hypertarget{sourcepointer-3}{}\fi
\noindent\textit{Source reference.} Original lemma 06YF: sources/stacks-chow/tags.txt, source lines 8884--8884. Cross-chapter standard prerequisite cited by the author inside the counted span; it is used as a pointer and its own proof is not part of this item.\par\medskip
\ifdefined\hypertarget\hypertarget{sourcepointer-4}{}\fi
\noindent\textit{Source reference.} Original lemma 0GT5: sources/stacks-chow/tags.txt, source lines 20955--20955. Cross-chapter standard prerequisite cited by the author inside the counted span; it is used as a pointer and its own proof is not part of this item.\par\medskip
\ifdefined\hypertarget\hypertarget{sourcepointer-5}{}\fi
\noindent\textit{Source reference.} Original lemma 0F8P: sources/stacks-chow/tags.txt, source lines 19049--19049. Cross-chapter standard prerequisite cited by the author inside the counted span; it is used as a pointer and its own proof is not part of this item.\par\medskip
\ifdefined\hypertarget\hypertarget{sourcepointer-6}{}\fi
\noindent\textit{Source reference.} Original reference 04VG: sources/stacks-chow/tags.txt, source lines 6324--6324. Cross-chapter standard prerequisite cited by the author inside the counted span; it is used as a pointer and its own proof is not part of this item.\par\medskip
\begin{lemma}
\label{lemma-Leta}
In Situation \ref{situation-eta} there is an additive
functor\footnote{Beware that this functor isn't exact, i.e.,
does not transform distinguished triangles into distinguished triangles.}
$L\eta_\mathcal{I} : D(\mathcal{O}_X) \to D(\mathcal{O}_X)$
such that if $M$ in $D(\mathcal{O}_X)$ is represented by a complex
$\mathcal{F}^\bullet$ of $\mathcal{I}$-torsion free $\mathcal{O}_X$-modules,
then $L\eta_\mathcal{I}M = \eta_\mathcal{I}\mathcal{F}^\bullet$.
Similarly for morphisms.
\end{lemma}

\begin{proof}
Denote $\mathcal{T} \subset \textit{Mod}(\mathcal{O}_X)$ the full subcategory
of $\mathcal{I}$-torsion free $\mathcal{O}_X$-modules.
We have a corresponding inclusion
$$
K(\mathcal{T})
\quad\subset\quad
K(\textit{Mod}(\mathcal{O}_X)) = K(\mathcal{O}_X)
$$
of $K(\mathcal{T})$ as a full triangulated subcategory of $K(\mathcal{O}_X)$.
Let $S \subset \text{Arrows}(K(\mathcal{T}))$ be the quasi-isomorphisms.
We will apply
Derived Categories, Lemma \ref{derived-lemma-localization-subcategory}
to show that the map
$$
S^{-1}K(\mathcal{T}) \longrightarrow D(\mathcal{O}_X)
$$
is an equivalence of triangulated categories. The lemma shows that
it suffices to prove: given a complex $\mathcal{G}^\bullet$ of
$\mathcal{O}_X$-modules,
there exists a quasi-isomorphism $\mathcal{F}^\bullet \to \mathcal{G}^\bullet$
with $\mathcal{F}^\bullet$
a complex of $\mathcal{I}$-torsion free $\mathcal{O}_X$-modules.
By Lemma \ref{lemma-K-flat-resolution} we can find a quasi-isomorphism
$\mathcal{F}^\bullet \to \mathcal{G}^\bullet$ such that the complex
$\mathcal{F}^\bullet$ is K-flat (we won't use this) and
consists of flat $\mathcal{O}_X$-modules $\mathcal{F}^i$.
By the third characterization of Lemma \ref{lemma-I-torsion-free}
we see that a flat $\mathcal{O}_X$-module is an
$\mathcal{I}$-torsion free $\mathcal{O}_X$-module
and we are done.

\medskip\noindent
With these preliminaries out of the way we can define $L\eta_f$.
Namely, by the discussion following Lemma \ref{lemma-I-torsion-free}
this section we have already a well defined functor
$$
K(\mathcal{T}) \xrightarrow{\eta_f} K(\mathcal{T}) \to
K(\mathcal{O}_X) \to D(\mathcal{O}_X)
$$
which according to Lemma \ref{lemma-eta-qis} sends quasi-isomorphisms
to quasi-isomorphisms. Hence this functor factors over
$S^{-1}K(\mathcal{T}) = D(\mathcal{O}_X)$ by 
Categories, Lemma \ref{categories-lemma-properties-left-localization}.
\end{proof}

\end{document}
